\section{Compressed direct proof-certificate quartics}
\label{sec:round13-direct}

The final selected direct construction gives a literal quartic in 1,859
nonnegative variables and a literal quartic in 3,179 unrestricted integer
variables.  The natural equation has 1,581 quadratic residuals, 12,275
monomials and coefficient height 63,924,066.  The integer equation has
2,074 residuals, 23,157 monomials and height 63,923,766.  Both heights are
below $2^{26}$.  For its integer polynomial $Q_8$, the theorem is
\[
 \operatorname{Con}(\mathrm{ZFC})
 \quad\Longleftrightarrow\quad
 \neg\exists z\in\mathbb Z^{3179}\;Q_8(z)=0.
\]
This is a no-zero equivalence, not a proof that the polynomial has no zero,
and not a new kernel-checked end-to-end translation.

We first establish the arithmetic and an eleven-column intermediate
source, then project to eight columns and normalize auxiliary signs in
Section~\ref{sec:projected-direct}.  The reconstructible intermediate
checkpoint has 2,170 natural variables or 4,087 integer variables; these
are not the final selected counts.  A retained sparse alternative with
2,560 natural variables and 11,074 monomials is also supplied.  Its smaller
monomial count is a genuine tradeoff against the selected 1,859-variable
natural polynomial.  The already improved baseline before this round has
3,443 natural variables, 2,838 residuals, 15,754 monomials and coefficient
height 127,302,193.  Late incoming integer comparisons and their advantages
are stated at the end of this section.

The construction follows the direct numerical C222 proof calculus,
independently of universal-polynomial or small-machine routes.  Its
calibration and indexed-Pell theorem are explicit inherited dependencies.
All new transformations below are existence-preserving mathematical
arguments, and the final files contain literal integer coefficients.

\subsection{The exact direct source and the carrier argument}

Write $I(x,y)=2(x+y)(x+y+1)+4y$, and define
\[
 A(x,y)=I(x,y)+2,\quad M(x,y)=I(x,y)+1,\quad
 E(x,y)=I(x,y)+3,\quad \neg x=I(x,1).
\]
These are numerical constructors, not arithmetic assertions about sets.
The four schema coordinates $a,b,c,d$ and twelve padded Horn branches are
exactly the pinned numerical C222 source.  The target is code $1$ and the
free truth code is $0$.  The explicit axiom constructors, including the
bundled axiom and its disequality guard, occur literally in
\texttt{build\_compact.py}; the inherited sparse source and calibration
modules are retained in \texttt{sources/}.

The dependency is the C222-to-ZF refutation theorem followed by the
constructible-choice ZF-to-ZFC refutation theorem, not execution of a
small universal machine.  In particular, no assumption that every natural
number decodes a well-formed formula is silently added.

For $N,m\ge1$, let $M=(m+1)^N$ and bound each rule label $p_i,q_i,v_i$
by $m$.  The accumulator source is
\[
 S_0=M+1,\quad S_N=1,\qquad
 S_{i+1}(1+v_iM)=S_i(1+p_iM)(1+q_iM),\quad i<N.
 \tag{1}
\]
Every row also requires the pinned Horn rule $(p_i,q_i)\to v_i$.
Telescoping gives an equality at $M$ between
\[
 F(X)=\prod_{i<N}(1+v_iX),\qquad
 G(X)=(1+X)\prod_{i<N}(1+p_iX)(1+q_iX).
\]
If $\deg G\ge N+1$, nonnegative coefficients give
$G(M)\ge M^{N+1}$, whereas
\[
 F(M)\le(1+mM)^N<((m+1)M)^N=M^{N+1}.
\]
Thus both products have at most $N$ nonconstant factors.  Their coefficient
sums at $1$ are at most $(m+1)^N=M$, and every individual coefficient is
strictly below $M$: a nonconstant product has its positive constant term
and another positive coefficient, while the constant product is $1<M$.
A nonzero polynomial $F-G$ cannot vanish at $M$, since its least nonzero
coefficient would be a nonzero multiple of $M$ with absolute value below
$M$.  Hence $F=G$.

The distinct factors $1+jX$, $j>0$, now recover the exact nonnegative
Horn flow balance for target $1$.  A finite single-conclusion flow implies
derivability: if the target were not derivable, sum the flow coordinates
outside the deductive closure.  Each valid rule contributes at most zero,
but the target contributes one.  Conversely, expand a finite proof tree
from the open obligation $1$ and let $S_i$ be the product of the weights
of all open obligations.  This gives (1).  Truth padding has weight one.
The proof uses single conclusions and nonnegative multiplicities; it does
not justify arbitrary cyclic proofs or signed flows.

\subsection{Eleven coordinates instead of thirteen}

Use the beta sequence function
\[
 \beta(x,y,i)=x\bmod\bigl(1+(i+1)y\bigr).
\]
For every finite natural sequence, its scale $y$ may be chosen strictly
larger than every sequence entry and divisible by $N!$.  The moduli for
$i=0,\ldots,N$ are then pairwise coprime: a common divisor is coprime to
$y$ and divides a positive index difference at most $N$, which divides
$y$.  The Chinese remainder theorem supplies $x$.

Parameterize
\[
 N=N_0+1,\quad m=m_0+1,\quad M=M_0+1=(m+1)^N,
 \qquad y=M_0+2+y_0,
\]
with all displayed underlying parameters natural.  The two endpoint rows
are simply
\[
 x=M_0+2+(y+1)u_0,
 \qquad x=1+\bigl(1+(N+1)y\bigr)u_N. \tag{2}
\]
They are canonical remainder statements because $M_0+2\le y<y+1$ and
$1<1+(N+1)y$.  The parameterization alone permits $S_0=y$; completeness
chooses $y>S_0$, and the first bounded row also enforces this strict
inequality.  No incorrect strict endpoint inequality is assumed.

At row $i$, retain only
\[
 w=(p,q,v,a,b,c,d,S,S',u,u')\in\mathbb N^{11}.
\]
Besides the Horn condition and the accumulator row, require
\[
 x=S+\bigl(1+(i+1)y\bigr)u,
 \qquad x=S'+\bigl(1+(i+2)y\bigr)u'. \tag{3}
\]
The CRT compiler bounds $p,q,v$ by $m$, bounds $S,S'$ by $y-1$, and
bounds $a,b,c,d,u,u'$ by one independent $U$.  Therefore (3) already
forces the canonical beta digits.  The two local remainder slack
coordinates and their two binomial range certificates disappear.

For completeness first choose the proof, $N,m,M$, and its finite
accumulator sequence, then choose $y$ strictly above every state, then
choose $x$ by CRT, and finally choose $U$ above the finitely many schema
and quotient coordinates.  None of these choices depends on the later
bounded-universal CRT encoding.  This order excludes circular bounds.

\subsection{Factoring the common Horn conditions}

Let $f_1,\ldots,f_8$ be the eight unconditional axiom codes, let $b_{\rm bundle}$ denote
the bundled axiom code, and let
\[
 g=((c-a)(c-b))^2-(d+1)^2
\]
be its guard residual.  Here $b$ in the guard remains the schema coordinate.
Put
\[
 J=v\prod_{r=1}^{8}(v-f_r),\qquad
 K=(v-b_{\rm bundle})^2+g^2.
\]
The following fixed nonnegative polynomial has exactly the pinned Horn
zero set:
\[
 H_{\rm Horn}=(q-I(p,v))^2
 \left[q^2+(v-A(a,p))^2(p^2+J^2)(p^2+K)\right]. \tag{4}
\]
Indeed its first factor is modus ponens.  Otherwise the bracket vanishes
exactly when $q=0$ and at least one of these alternatives holds:
$v=A(a,p)$; or $p=0$ and $J=0$; or $p=0$ and $K=0$.
These are generalization, truth or one of the eight unconditional axioms,
and the guarded bundled axiom, respectively.  Every factor used in this
argument is nonnegative.  This is an equivalence of zero sets, not an
asserted polynomial identity with the old twelve-factor source.

Let $P$ be (4) plus the squares of the accumulator residual and the two
residuals in (3).  The signed arithmetic circuit is split as $P=E-F$,
where $E,F$ have nonnegative coefficients, using
\[
 (e,f)+(g,h)=(e+g,f+h),\qquad
 (e,f)(g,h)=(eg+fh,eh+fg),
\]
and $(e-f)^2=(e^2+f^2)-2ef$.  Thus $E=F$ holds exactly when every
required local condition holds.

\subsection{One small circuit for a very large majorant}

It is unnecessary to evaluate a second full copy of the source circuit at
all coordinate upper bounds.  Put
\[
 R=4+N+x+y+m+M_0+U.
\]
Every input at which the upper source is evaluated is at most $R$: the
index is $N$, the label bounds are $m$, the six other private bounds are
$U$, and $y$ is an upper bound for each state (slightly larger than its
actual bound $y-1$).  The external inputs are also included in $R$.

A structural exponent certificate bounds every nonnegative circuit node
by a power of $R$.  An input receives exponent one.  A constant $c$
receives the least $e$ with $c\le4^e$.  At addition take
$\max(e_1,e_2)+1$, since $R\ge4\ge2$; at multiplication take $e_1+e_2$.
Induction on the circuit proves the bound at every node.  For the actual
factored source the root exponents are 6,893,154 and 6,893,148.
Rounding upwards to $2^{23}=8,388,608$ gives
\[
 H=R^{2^{23}}\ge N,\qquad H\ge E(N,z,u),F(N,z,u). \tag{5}
\]
Here the upper state entries in $u$ may be taken to be $y$, as just noted.
The literal compiler computes $H$ by exactly 23 successive squarings.
There is no variable exponent or power-representation macro at this step.
The larger numerical value of the bound has essentially no circuit cost.
The exported exponent certificate records every abstract circuit node
and every one of these 23 squarings.

\subsection{One bounded-quantifier certificate with direct cancellation}

Set
\[
 T=H!,\quad D_i=1+(i+1)T,\quad B=\prod_{i<N}D_i.
\]
For actual coordinate bounds
$u=(m,m,m,U,U,U,U,y-1,y-1,U,U)$, the bounded row family is equivalent to
existence of positive encoded columns $A_j$ and a natural encoded index $C$
with
\[
 B\mid {A_j\choose u_j+1},\qquad
 B\mid1+(C+1)T,\qquad E(C,z,A)\equiv F(C,z,A)\pmod B. \tag{6}
\]
In particular, the two state binomial lower arguments are exactly $y$;
no subtraction is needed in the finite arithmetic circuit.

For soundness fix $i<N$ and take a prime $p\mid D_i$.  It is coprime to
$T=H!$, hence $p>H$.  The identity
\[
 (u_j+1)!{A_j\choose u_j+1}=\prod_{r=0}^{u_j}(A_j-r)
\]
implies that $A_j\equiv w_j\pmod p$ for some $0\le w_j\le u_j$.
This remains true when $A_j\le u_j$, as one factor is zero.  The index
condition gives $C\equiv i\pmod p$.  The evaluation congruence then holds
for $E(i,z,w),F(i,z,w)$; both lie in $[0,H]$, so they are equal.

For completeness the $D_i$ are pairwise coprime because $H\ge N$.
The radius bound supplies more than needed for soundness: $H\ge R$ and
$R\ge u_j+1$ for every coordinate.  Consequently
$L_j=(u_j+1)!$ divides $H!=T$, and $\gcd(B,L_j)=1$.
Choose ordinary CRT representatives
$A_j\equiv w_j^{(i)}\pmod{D_i}$ and $C\equiv i\pmod{D_i}$.
The falling-factorial numerator then vanishes modulo every $D_i$ and
hence modulo $B$.  Cancelling $L_j$, which is coprime to $B$, gives
$B\mid {A_j\choose u_j+1}$.  Adding a sufficiently large multiple of
$B$ makes $A_j>0$, preserving the falling-product congruence and hence,
by the same cancellation, the binomial divisibility.  The valid rows
and polynomial congruence give the evaluation condition in (6).

This final radius construction therefore needs neither prime-supported
CRT lifts nor their proof-only products.  Those are necessary in the more
general current-baseline theorem without a factorial-size guard, but the
very cheap radius bound restores the simpler direct cancellation proof.
There is no circularity: all coordinate bounds and $H$ are chosen before
$T,B$ and the encoded columns.

\subsection{All arithmetic operators are eliminated}

The exact arithmetic identities are
\[
 Q_0=(1+(N+1)T)^N+2,\qquad
 B=\left(T^N N!{Q_0+N\choose N}\right)\bmod(TQ_0-1), \tag{7}
\]
\[
 {n\choose k}=
 \left\lfloor\frac{(2^n+2)^n}{(2^n+1)^k}\right\rfloor\bmod(2^n+1), \tag{8}
\]
\[
 r=(n+2)^{n+2},\quad C_0={r+n\choose n},\qquad
 n!=\left\lfloor\frac{(r+1)^n+C_0-1}{C_0}\right\rfloor. \tag{9}
\]
For (7), reduction of the numerator modulo $TQ_0-1$ gives
$\prod_{j=1}^N(1+jT)$, and the chosen $Q_0$ makes this product strictly
smaller than the modulus.  Division uses $a=dq+r$ and $d=r+s+1$.
The circuit for (8) combines its two divisions using one high-digit
coordinate.  Formula (9) uses the equivalent equation
$(r+1)^n+C_0=C_0 n!+r_0+1$, with $0\le r_0<C_0$.
The included factorial and bounded-forall representation sources state
these identities with their boundary cases.

There are exactly 49 positive-power graph calls, 14 binomial calls, and
two factorial calls.  Positivity is established before each call:
$N,m\ge1$; $H\ge4$; encoded columns are parameterized as $A_j=\alpha_j+1$;
range lower arguments are $m+1$, $U+1$, or $y\ge2$; all auxiliary
factorial/binomial upper arguments and power bases are positive.
Thus deleting zero-base and zero-exponent cases is licensed for this
construction, not for the unrestricted exponentiation relation.

\subsection{The parametric positive-power graph}

For positive $n,k$, let $m$ be the proposed output.  There are twelve
natural certificate coordinates
\[
 a_2,y_k,g,d,X,U,q_4,q_B,q_S,q_R,q_P,Z.
\]
Define, rather than allocate additional unconstrained coordinates,
\begin{align*}
 W&=n+k,& A&=a_2+2,& Y&=k+y_k,& T_0&=m+g+1,\\
 V&=Y^2(d+1),& B_0&=1+4Y(q_4+1),&
 S&=X+Uq_S,& R_0&=k+4Yq_R.
\end{align*}
Writing $\mathcal N(a,x,y)$ for
$x^2+y^2-a^2y^2=1$, require precisely seven polynomial conditions:
\begin{gather*}
 \mathcal N(A,X,Y),\quad \mathcal N(A,U,V),\quad
 \mathcal N(B_0,S,R_0),\\
 B_0=A+Uq_B,\qquad
 X+Yn=YA+m+T_0q_P,\qquad 2An=T_0+n^2+1,\\
 \mathcal N(W+1,A,WZ). \tag{10}
\end{gather*}
These are circuit expressions; flattening their products gives quadratic
residuals.  The explicitly assigned $B_0$ is strictly greater than one,
including when $q_4=0$.  There is no missing bound on this Pell parameter.

\paragraph{Soundness.}
The definitions give $A>1$, $Y\ge k>0$, $V>0$, $Y^2\mid V$, and
$B_0>1$.  The three forward congruence equations imply the three modular
conditions of the pinned indexed-Pell criterion.  Thus $(X,Y)$ is the
$k$th Pell solution for parameter $A$.  The last norm gives
\[
 A^2=1+W^3(W+2)Z^2.
\]
As $A\ge2$ and $W\ge1$, necessarily $Z\ge1$ and $A>W\ge n$.
Therefore the rearranged power congruence uses the same ordinary
subtraction as the pinned power theorem.  Also $m<T_0$, and the deleted
common-bound slacks can be restored as $W-n=k$ and $W-k=n$.
The positive-power criterion yields $m=n^k$.

\paragraph{Completeness and the direction of every quotient.}
Let $(X_a(j),Y_a(j))$ denote the natural Pell sequence.  At $W=n+k$ choose
\[
 A=X_{W+1}(W),\quad WZ=Y_{W+1}(W),\quad
 T_0=2An-n^2-1,\quad X=X_A(k),\ Y=Y_A(k).
\]
The standard congruence $Y_{W+1}(W)\equiv W\pmod W$ supplies $Z$.
The usual fixed-bound power construction gives $n^k<T_0$ and the power
congruence.  For its direction set
$F_j=X_A(j)-(A-n)Y_A(j)$.  Then
$F_0=1,F_1=n,F_{j+1}=2AF_j-F_{j-1}$.  Inductively
$F_j\ge F_{j-1}\ge0$ and
$F_{j+1}\ge(2A-1)F_j\ge nF_j$, so $F_k\ge n^k$.
Thus $q_P=(F_k-n^k)/T_0$ is natural.

In the explicit completeness proof of the indexed-Pell criterion set
$h=2kY$, $U=X_A(h)$, $V=Y_A(h)$, and choose $B_0$ by CRT with
$B_0\equiv A\pmod U$ and $B_0\equiv1\pmod{4Y}$.
That proof has $A<U$, so its natural CRT representative satisfies
$B_0\ge A$.  Take $S=X_{B_0}(k)$ and $R_0=Y_{B_0}(k)$, at the same
index $k$.  The Pell sequences are monotone in the parameter, hence
$S\ge X$; also $R_0\ge k$.  Consequently all $q_B,q_S,q_R$ are natural.
Because $B_0>1$ and $Y>0$, its quotient in $B_0=1+4Yq$ satisfies $q\ge1$,
so write $q=q_4+1$.  The remaining slacks follow from the displayed
inequalities and divisibility.  This proves completeness of (10), not
merely its one-way soundness.  The separate orientation supplement
records the monotonicity proof and exact pinned source lines.

\subsection{Quartic flattening, sparse alternative, and integer domain}

Every nonnegative arithmetic addition or multiplication receives a
triangular defining row.  Nonnegative affine gate values may be
substituted without raising residual degree; the implementation substitutes
only expressions of support at most four, a size tradeoff rather than a
semantic restriction.  The complete substitution map is supplied.
The final polynomial is the sum of squares of all remaining quadratic
rows.  Over either $\mathbb N$ or $\mathbb Z$, a sum of integer squares
vanishes exactly when all rows vanish.  Exact expansion and collection
give the reported quartics.

The 2,560-variable sparse alternative retains semantic names for the Pell
outputs rather than making all of (10)'s definitions affine substitutions.
It has fewer collected monomials.  Its orientation, short-beta source,
CRT block and radius bound are otherwise the same.  The intermediate
shared-offset construction is also justified: finitely many signed
quotients can be written $v_i-h$ with one common natural $h$, chosen after
all other witnesses.  It is superseded by the forward form for the final
natural variants, except that one final CRT congruence still uses two
natural quotients.

For the 4,087-variable integer companion, congruence quotients themselves
need not be nonnegative.  There are exactly 197 unrestricted quotient
coordinates: four per positive-power graph and one final CRT quotient.
In (10), $S,R_0$ are retained as separate natural semantic coordinates,
with their defining equalities, while $q_S,q_R$ may be signed.  The other
parametric definitions are unchanged.  In particular $B_0>1$ still follows
from $B_0=1+4Y(q_4+1)$ with natural $q_4$; it is not incorrectly inferred
from signed $q_B$.  Soundness reads congruences directly and applies the
same indexed-Pell and power theorems.  Completeness uses the natural
witness construction already proved.  The final CRT quotient may be
signed because (6) is a divisibility condition.

Only 607 further nonnegativity assertions are necessary.  An exact sign
certificate infers the other 44 of the 651 nonquotient semantic
coordinates from polynomial equalities.  Each inference is of the form
$c x=P$ with $c>0$ and $P$ having nonnegative coefficients; every variable
occurring to an odd exponent in a monomial of $P$ must already be known
nonnegative.  Even powers need no sign premise.  The supplied application
order is an acyclic proof of all required signs.  Signed quotient gates
are not asserted to be nonnegative.

For each remaining seed $x$, append the quadratic residual
\[
 2x-a(a+1)-b(b+1)-c(c+1)=0. \tag{11}
\]
For arbitrary integers $a,b,c$, each consecutive product is nonnegative,
so (11) forces $x\ge0$.  Conversely every natural $x$ is a sum of three
triangular numbers by Gauss's theorem.  Equivalently, apply Legendre's
three-square theorem to $8x+3$: the three squares must all be odd modulo
8, and subtracting their three units yields (11).
The 607 seeds add 1,821 integer coordinates and 607 quadratic rows to
the 2,266-coordinate mixed-domain template.  This proves the counts
$2266+1821=4087$ and $1924+607=2531$, and preserves degree four.

\subsection{Source pins, reproduction, and the claim boundary}

The calibration files are pinned to
\texttt{VladimirReshetnikov/ProveIt} commit
\nolinkurl{08200eb10fde2db408c69dffd840d6920a569b1f};
the comparison source is the fixed-power-bound candidate at commit
\nolinkurl{7e533c33819ee1f703b69f2778a767cdc069940d}.
The exact theorem names, file paths and SHA256 values are listed in
\texttt{SOURCE\_PINS.json} and the inherited calibration inventory.
The all-forward Pell proof additionally pins the official Mathlib file
at commit \nolinkurl{4dc43842a11f257ffba2a611c943a133a65cb4b2},
SHA256 \nolinkurl{5c5756cdf2fa70fb77af9daa9c7167c65749ac616bde8f5fdfeed2cb958b3ed5}.
It is inspected as source text, not executed.

The builders are newly authored standard-library Python implementations.
No upstream executable was imported or run.  Their baseline mode
independently reconstructs all five published numerical baseline metrics.
The final release includes the exact sparse equations, every residual,
raw circuits and macro bindings, affine restoration maps, structural
radius certificates, integer sign certificates, source pins, and
independent finite-audit receipts.  Reproduction commands are in
\texttt{README.md} and the file whitelist is machine-readable.

These finite checks establish literal polynomial and certificate facts,
not consistency of ZFC and not the universal mathematical dependencies.
The logical chain is: pinned numerical refutation calibration; finite
single-conclusion proof flow and the carrier lemma; short-beta source;
factored Horn zero set; bounded CRT with factorial cancellation; radius bound; explicit arithmetic
identities and Pell graphs; quartic flattening; and, for unrestricted
integers, sign closure and the three-triangular-number theorem.  No
nonpolynomial operation remains in the final equation.

\subsection{Final source projection: eight columns and no private bound}
\label{sec:projected-direct}

The source projection in this subsection is explicitly credited to the
incoming \emph{ConZFC Direct Projection 7617} package, whose preprocessing
proof is preserved as \texttt{sources/INCOMING\_PROJECTION\_PROOF.md}.
The new result here is its composition with the sharper carrier,
parametric positive-power graphs, dyadic majorants, quotient bound $x$,
and sign-normalized integer compilation.  We do not claim that eliminating
these three label columns was first discovered in this construction.

Each rule can be parameterized by just $a,b,c,d$:
\[
\begin{array}{c|c|c|c}
 \text{rule}&p&q&v\\ \hline
 \text{truth}&0&0&0\\
 \text{eight ordinary axioms}&0&0&f_j(a,b,c,d)\\
 \text{bundled axiom}&0&0&f_{\rm bundle}(a,b,c,d)\\
 \text{modus ponens}&a&I(a,b)&b\\
 \text{generalization}&b&0&A(a,b)
\end{array}
\]
The bundled row additionally requires
$((c-a)(c-b))^2=(d+1)^2$.  No other row has a guard.  These substitutions
are exactly the pinned C222 rules, including the independent
parameter and premise of generalization.

For a chosen branch $j$, substitute its labels $(p_j,q_j,v_j)$ into the
accumulator residual
\[
 L_j=S'(1+v_jM)-S(1+p_jM)(1+q_jM).
\]
Let $C_j=L_j^2$ except at the bundle, where
$C_j=L_j^2+g^2$.  The new local polynomial is
\[
 P_8=\prod_{j=0}^{11}C_j
 +\bigl(x-S-(1+(i+1)y)u\bigr)^2
 +\bigl(x-S'-(1+(i+2)y)u'\bigr)^2. \tag{12}
\]
Its zero set is exactly one valid branch-specific update together with
the two common beta reconstructions.  As all branch factors are
nonnegative, the product implements the finite disjunction without
selecting a rule by an extra coordinate.  This is existential projection
of the three deterministic labels, not an assertion that the polynomial
in (12) is algebraically identical to the eleven-coordinate source.

\paragraph{Padding chooses the parameter bound before the carrier.}
Given a finite proof-tree expansion, choose $N\ge1$ at least its length
and at least every schema parameter in the expansion.  For modus ponens,
this may include a very large formula code; no practical bound on the
resulting $N$ is claimed.  Pad to exactly $N$ rows using truth, whose
labels are $(0,0,0)$ and whose update leaves $S$ unchanged.  Thus all four
parameters may be bounded by $N$ without changing refutability.

Every branch label is a nonnegative-coefficient polynomial in those four
parameters.  Apply the structural exponent certificate to its literal
constructor DAG with radius $N+4$.  The largest certified exponent is
3,445,903, so the fixed polynomial
\[
 m=(N+4)^{2^{22}} \tag{13}
\]
bounds every label.  It is computed by exactly 22 squarings, not by a
variable-exponent graph and not by a duplicate evaluation of the bundled
axiom.  It satisfies $m\ge1$.  Choose $M=(m+1)^N$ and $M_0=M-1$.
The carrier proof already given applies verbatim to these labels.  The
choice order is proof, padded $N$, then $m,M$, then the accumulator
sequence, then $y,x$.  There is no self-referential bound.

The eight local coordinates are
\[
 (a,b,c,d,S,S',u,u'),\qquad
 \text{bounds }(N,N,N,N,y-1,y-1,x,x). \tag{14}
\]
The use of $x$ for the two quotient bounds is an additional simplification:
canonical beta reconstructions have
$x=S+Du$ with $S\ge0,D\ge1$, hence $u\le x$, and similarly for $u'$.
No new private bound $U$ is needed.  Conversely, the state bounds in (14)
make those reconstructions canonical, just as in the eleven-column proof.
The endpoints and parameterization $y=M_0+2+y_0$ are unchanged.
The four outer natural coordinates are $N_0,x,y_0,M_0$, with $N=N_0+1$.

\paragraph{A second dyadic radius proves the shared CRT bound.}
Split the actual signed circuit for (12) as $P_8=E_8-F_8$ and put
\[
 R=4+N+x+y+M_0,\qquad H=R^{2^{23}}. \tag{15}
\]
All inputs at the upper substitution of (14) are at most $R$; for the two
states we safely use $y$ instead of $y-1$.  The two root exponents from
the structural certificate are both 6,893,248, below $2^{23}$.
Therefore $H$ bounds both source evaluations and $N$.  It also dominates
every binomial cutoff $N+1,y,x+1$.  The same CRT theorem, with direct
factorial cancellation, applies.  Only 23 successive squarings compute
(15).  Notice that $m$ need not occur in this source radius: after label
projection the local source computes its labels from the four bounded
parameters and uses only $M_0$ as its carrier input.

The final eight-column arithmetic uses exactly 40 positive-power graphs,
11 binomial graphs and two factorial graphs.  The count is transparent:
24 powers for the eight coordinate binomials, five for each of the two
factorials, five other powers in the modulus product, and one carrier
power.  In this accounting the five other product powers are its two
explicit powers and the three powers of its extra binomial; the factorial
at $N$ has already been counted.  No arithmetic macro, variable-length
product or variable exponent remains after flattening.

\subsection{Existential sign normalization for the final integer equation}

The incoming signed-quartic package supplies the sign-orbit method
and its two pivots per positive-power graph; its proof is preserved as
\texttt{sources/INCOMING\_SIGN\_NORMALIZATION\_PROOF.md}.  The final
application here is checked against the new projected parametric system,
not assumed merely because it worked for a predecessor.

For a family of residuals $r_j$, a diagonal sign action
$x_i\mapsto(-1)^{\epsilon_i}x_i$ preserves the common zero set whenever
all monomials of each row have equal parity against $\epsilon$.
Then $r_j$ is changed to either $r_j$ or $-r_j$.
For each of the 40 power calls there are two independent actions:
\begin{itemize}
\item Flip its auxiliary $Z$ and the product wire $WZ$.  The only other
      use of that wire is through its square in the last norm.
\item Flip its Pell coordinate $U$ and its two unrestricted quotients
      $q_B,q_S$.  The products $Uq_B,Uq_S$ are unchanged, while the norm
      uses $U$ only through $U^2$.
\end{itemize}
The actions leave every other semantic coordinate unchanged.  They
commute, and do not change any retained sign guard.  At any integer zero,
apply a pivot's action exactly when that pivot is negative.  This produces
another zero in which all 80 pivots are nonnegative.  This is an
existence-preserving normalization, not a claim that the original zero
already had those signs.

The released normalization certificate specifies all flipped raw and
exported variables and the parity of every residual under each action.
The separate independent checker verifies these finite identities; the
argument above explains their universal meaning.  Gate values involving
signed quotients remain allowed to be negative.  No signs of $X,S,R_0$
are discarded by an unsupported symmetry argument.

The mixed-domain template has 1,937 variables and 1,660 quadratic rows.
It has 161 unrestricted congruence quotients: four per power call and
one final CRT quotient.  There are 529 other semantic coordinates.
Before normalization, exact polynomial sign closure requires 494 seeds
and infers 35 more signs.  The 80 pivots are among these seeds and can
instead be normalized.  Thus only
\[
 414+80+35=529
\]
semantic signs need to be guarded, normalized, or inferred, respectively.
The 414 three-triangular-number certificates add $3\cdot414=1242$
integer variables and 414 quadratic rows.  Consequently the final
integer counts are
\[
 1937+1242=3179,\qquad 1660+414=2074.
\]
The resulting ordinary quartic $Q_8$ has 23,157 monomials and coefficient
height 63,923,766.  The natural companion has 1,859 variables, 1,581 rows,
12,275 monomials and height 63,924,066.  Both heights are below $2^{26}$.
The literal files are \texttt{projected/expanded\_polynomial.json} and
\texttt{projected\_integer/expanded\_polynomial.json}.

\paragraph{Complete equivalence and comparisons.}
A zero of the integer polynomial first yields an integer zero of the
mixed template with the 414 guarded signs.  Normalize the 80 pivots and
replay the 35 sign implications.  All nonquotient arithmetic inputs are
then natural.  The parametric Pell certificates establish every required
power; formulas (7)--(9) establish every other arithmetic macro; the CRT
certificate recovers all eight-column rows; their parameter and label
bounds license the sharp carrier and hence a C222 refutation.  The pinned
calibration gives a ZFC refutation.  Conversely a finite refutation gives
the padded source in (12)--(14), all arithmetic certificates, natural
power witnesses, and the 414 triangular representations, hence an integer
zero.  Therefore
\[
 \operatorname{Con}(\mathrm{ZFC})
 \quad\Longleftrightarrow\quad
 \neg\exists z\in\mathbb Z^{3179}\;Q_8(z)=0.
\]
All steps are ordinary mathematics formalizable in ZFC.  Finite artifact
checks do not replace the stated unbounded arithmetic and calibration
theorems, and no new proof-assistant certification or proof of consistency
is claimed.

The late incoming direct integer alternatives have 3,776 variables,
20,194 monomials and height 47,959; and 3,671 variables, 26,271 monomials
and height 3,835,369,234.  The new 3,179-variable equation improves the
latter on variables, monomials and coefficient height, but does not
improve the former's monomial count or its much smaller coefficients.
It should not be described as dominating every current direct equation.
The 2,560-natural-variable sparse companion likewise remains relevant
because its 11,074 monomials are fewer than the projected natural
companion's 12,275.  These are explicit tradeoffs, not an absolute
minimality claim.
